%% ============================================================================
%%  MASTER RESUME  --  the single source of truth.
%%
%%  Everything you have ever done lives in this file. Two things read it:
%%    1. LaTeX  -> builds the full "master" PDF.
%%    2. The website (src/) -> parses this file at build time and renders
%%       Education, Skills, Experience, Projects and Extracurriculars.
%%
%%  So: edit here, and both the PDF and the website update. Never edit the
%%  same fact in two places.
%%
%%  ---- RULES SO THE WEBSITE CAN READ THIS FILE ------------------------------
%%  * Keep the section names:  Education, Technical Skills, Experience,
%%    Projects, Extracurriculars.
%%  * Experience heading:   {\textbf{Role - Project - Organisation}}
%%      (separate the parts with " - " = space, hyphen, space).
%%      The LAST part is shown as the organisation, the rest as the title.
%%  * Tech line:   {\emph{\textbf{Tech:}} A, B, C}   (comma separated)
%%  * Link line (projects only):
%%      {\\ \small \textbf{\emph{Link:}} \href{https://url}{text}}
%%  * One \resumeItem{...} per bullet point.
%%  * Skills:  \textbf{Category}{: item, item, item} \\
%%  * Lines starting with % are ignored by BOTH LaTeX and the website, so
%%    commenting an entry out hides it everywhere.
%%  * Contact details (phone, email, GitHub, portfolio) are NOT here; they
%%    live in ../shared/profile.tex.
%%
%%  Build the PDF:   make master-resume
%% ============================================================================

\documentclass[letterpaper,11pt]{article}
%% ============================================================================
%%  RESUME PREAMBLE  --  the look of every resume (master, temporary, final).
%%
%%  Edit the styling here once and all three resumes pick it up.
%%  Each resume file starts with:
%%      \documentclass[letterpaper,11pt]{article}
%%      %% ============================================================================
%%  RESUME PREAMBLE  --  the look of every resume (master, temporary, final).
%%
%%  Edit the styling here once and all three resumes pick it up.
%%  Each resume file starts with:
%%      \documentclass[letterpaper,11pt]{article}
%%      %% ============================================================================
%%  RESUME PREAMBLE  --  the look of every resume (master, temporary, final).
%%
%%  Edit the styling here once and all three resumes pick it up.
%%  Each resume file starts with:
%%      \documentclass[letterpaper,11pt]{article}
%%      \input{../shared/resume_preamble}
%% ============================================================================

\usepackage{latexsym}
\usepackage[empty]{fullpage}
\usepackage{titlesec}
\usepackage{marvosym}
\usepackage[usenames,dvipsnames]{color}
\usepackage{verbatim}
\usepackage{enumitem}
\usepackage[hidelinks]{hyperref}
\usepackage{fancyhdr}
\usepackage[english]{babel}
\usepackage{tabularx}
\input{glyphtounicode}

\input{../shared/profile}

%----------FONT OPTIONS----------

\pagestyle{fancy}
\fancyhf{} % clear all header and footer fields
\fancyfoot{}
\renewcommand{\headrulewidth}{0pt}
\renewcommand{\footrulewidth}{0pt}

% Adjust margins
\addtolength{\oddsidemargin}{-0.5in}
\addtolength{\evensidemargin}{-0.5in}
\addtolength{\textwidth}{1in}
\addtolength{\topmargin}{-.5in}
\addtolength{\textheight}{1.0in}

\urlstyle{same}

\raggedbottom
\raggedright
\setlength{\tabcolsep}{0in}

% Sections formatting
\titleformat{\section}{
  \vspace{-4pt}\scshape\raggedright\large
}{}{0em}{}[\color{black}\titlerule \vspace{-5pt}]

% Ensure that generate pdf is machine readable/ATS parsable
\pdfgentounicode=1

%-------------------------
% Custom commands
\newcommand{\resumeItem}[1]{
  \item\small{
    {#1 \vspace{-2pt}}
  }
}

\newcommand{\resumeSubheading}[4]{
  \vspace{0pt}\item
    \begin{tabular*}{0.97\textwidth}[t]{l@{\extracolsep{\fill}}r}
      \textbf{#1} & #2 \\
      \textit{\small#3} %& \textit{\small #4} \\
    \end{tabular*} % \vspace{-7pt}
    \textit{\small #4}
    \vspace{-7pt}
}

\newcommand{\skillSubheading}[4]{
  \vspace{1pt}\item
    \begin{tabular*}{0.97\textwidth}[t]{l@{\extracolsep{\fill}}r}
      %\textbf{#1} & #2 \\
      \textit{\small#3} & \textit{\small #4} \\
    \end{tabular*}\vspace{-7pt}
}

\newcommand{\resumeSubSubheading}[2]{
    \item
    \begin{tabular*}{0.97\textwidth}{l@{\extracolsep{\fill}}r}
      \textit{\small#1} & \textit{\small #2} \\
    \end{tabular*}\vspace{-7pt}
}

\newcommand{\resumeProjectHeading}[2]{
    \item
    \begin{tabular*}{0.97\textwidth}{l@{\extracolsep{\fill}}r}
      \small#1 & #2 \\
    \end{tabular*}\vspace{-7pt}
}

\newcommand{\resumeSubItem}[1]{\resumeItem{#1}\vspace{-4pt}}

\renewcommand\labelitemii{$\vcenter{\hbox{\tiny$\bullet$}}$}

\newcommand{\resumeSubHeadingListStart}{\begin{itemize}[leftmargin=0.10in, label={}]}
\newcommand{\resumeSubHeadingListEnd}{\end{itemize}}
\newcommand{\resumeItemListStart}{\begin{itemize}[leftmargin=0.14in]}
\newcommand{\resumeItemListEnd}{\end{itemize}\vspace{-8pt}}

%-------------------------
% Page header (name + contact line). Values come from shared/profile.tex.
\newcommand{\resumeHeader}{%
  \begin{center}
    \textbf{\Huge \scshape \profileName} \\ \vspace{5pt}
    \profilePhone\ $|$
    \href{mailto:\profileEmail}{\profileEmail} $|$
    \href{https://\profileGithub}{\profileGithub} $|$
    Portfolio: \href{https://\profilePortfolio}{\profilePortfolio}
  \end{center}%
}

%% ============================================================================

\usepackage{latexsym}
\usepackage[empty]{fullpage}
\usepackage{titlesec}
\usepackage{marvosym}
\usepackage[usenames,dvipsnames]{color}
\usepackage{verbatim}
\usepackage{enumitem}
\usepackage[hidelinks]{hyperref}
\usepackage{fancyhdr}
\usepackage[english]{babel}
\usepackage{tabularx}
\input{glyphtounicode}

%% ============================================================================
%%  PROFILE  --  your contact details, in ONE place.
%%
%%  Used by:  every resume, every cover letter, AND the website.
%%  Change a value here and it changes everywhere on the next build.
%%
%%  Keep each value on one line, written exactly as  \newcommand{\profileXxx}{value}
%%  (the website reads this file with a simple parser).
%%  Links are written WITHOUT "https://".
%% ============================================================================

\newcommand{\profileName}{Sharzil Hasnine}
\newcommand{\profilePhone}{+61 451 645 516}
\newcommand{\profileEmail}{sharzilhasnine975@gmail.com}
\newcommand{\profileGithub}{github.com/Shas-9}
\newcommand{\profilePortfolio}{shas-9.github.io}


%----------FONT OPTIONS----------

\pagestyle{fancy}
\fancyhf{} % clear all header and footer fields
\fancyfoot{}
\renewcommand{\headrulewidth}{0pt}
\renewcommand{\footrulewidth}{0pt}

% Adjust margins
\addtolength{\oddsidemargin}{-0.5in}
\addtolength{\evensidemargin}{-0.5in}
\addtolength{\textwidth}{1in}
\addtolength{\topmargin}{-.5in}
\addtolength{\textheight}{1.0in}

\urlstyle{same}

\raggedbottom
\raggedright
\setlength{\tabcolsep}{0in}

% Sections formatting
\titleformat{\section}{
  \vspace{-4pt}\scshape\raggedright\large
}{}{0em}{}[\color{black}\titlerule \vspace{-5pt}]

% Ensure that generate pdf is machine readable/ATS parsable
\pdfgentounicode=1

%-------------------------
% Custom commands
\newcommand{\resumeItem}[1]{
  \item\small{
    {#1 \vspace{-2pt}}
  }
}

\newcommand{\resumeSubheading}[4]{
  \vspace{0pt}\item
    \begin{tabular*}{0.97\textwidth}[t]{l@{\extracolsep{\fill}}r}
      \textbf{#1} & #2 \\
      \textit{\small#3} %& \textit{\small #4} \\
    \end{tabular*} % \vspace{-7pt}
    \textit{\small #4}
    \vspace{-7pt}
}

\newcommand{\skillSubheading}[4]{
  \vspace{1pt}\item
    \begin{tabular*}{0.97\textwidth}[t]{l@{\extracolsep{\fill}}r}
      %\textbf{#1} & #2 \\
      \textit{\small#3} & \textit{\small #4} \\
    \end{tabular*}\vspace{-7pt}
}

\newcommand{\resumeSubSubheading}[2]{
    \item
    \begin{tabular*}{0.97\textwidth}{l@{\extracolsep{\fill}}r}
      \textit{\small#1} & \textit{\small #2} \\
    \end{tabular*}\vspace{-7pt}
}

\newcommand{\resumeProjectHeading}[2]{
    \item
    \begin{tabular*}{0.97\textwidth}{l@{\extracolsep{\fill}}r}
      \small#1 & #2 \\
    \end{tabular*}\vspace{-7pt}
}

\newcommand{\resumeSubItem}[1]{\resumeItem{#1}\vspace{-4pt}}

\renewcommand\labelitemii{$\vcenter{\hbox{\tiny$\bullet$}}$}

\newcommand{\resumeSubHeadingListStart}{\begin{itemize}[leftmargin=0.10in, label={}]}
\newcommand{\resumeSubHeadingListEnd}{\end{itemize}}
\newcommand{\resumeItemListStart}{\begin{itemize}[leftmargin=0.14in]}
\newcommand{\resumeItemListEnd}{\end{itemize}\vspace{-8pt}}

%-------------------------
% Page header (name + contact line). Values come from shared/profile.tex.
\newcommand{\resumeHeader}{%
  \begin{center}
    \textbf{\Huge \scshape \profileName} \\ \vspace{5pt}
    \profilePhone\ $|$
    \href{mailto:\profileEmail}{\profileEmail} $|$
    \href{https://\profileGithub}{\profileGithub} $|$
    Portfolio: \href{https://\profilePortfolio}{\profilePortfolio}
  \end{center}%
}

%% ============================================================================

\usepackage{latexsym}
\usepackage[empty]{fullpage}
\usepackage{titlesec}
\usepackage{marvosym}
\usepackage[usenames,dvipsnames]{color}
\usepackage{verbatim}
\usepackage{enumitem}
\usepackage[hidelinks]{hyperref}
\usepackage{fancyhdr}
\usepackage[english]{babel}
\usepackage{tabularx}
\input{glyphtounicode}

%% ============================================================================
%%  PROFILE  --  your contact details, in ONE place.
%%
%%  Used by:  every resume, every cover letter, AND the website.
%%  Change a value here and it changes everywhere on the next build.
%%
%%  Keep each value on one line, written exactly as  \newcommand{\profileXxx}{value}
%%  (the website reads this file with a simple parser).
%%  Links are written WITHOUT "https://".
%% ============================================================================

\newcommand{\profileName}{Sharzil Hasnine}
\newcommand{\profilePhone}{+61 451 645 516}
\newcommand{\profileEmail}{sharzilhasnine975@gmail.com}
\newcommand{\profileGithub}{github.com/Shas-9}
\newcommand{\profilePortfolio}{shas-9.github.io}


%----------FONT OPTIONS----------

\pagestyle{fancy}
\fancyhf{} % clear all header and footer fields
\fancyfoot{}
\renewcommand{\headrulewidth}{0pt}
\renewcommand{\footrulewidth}{0pt}

% Adjust margins
\addtolength{\oddsidemargin}{-0.5in}
\addtolength{\evensidemargin}{-0.5in}
\addtolength{\textwidth}{1in}
\addtolength{\topmargin}{-.5in}
\addtolength{\textheight}{1.0in}

\urlstyle{same}

\raggedbottom
\raggedright
\setlength{\tabcolsep}{0in}

% Sections formatting
\titleformat{\section}{
  \vspace{-4pt}\scshape\raggedright\large
}{}{0em}{}[\color{black}\titlerule \vspace{-5pt}]

% Ensure that generate pdf is machine readable/ATS parsable
\pdfgentounicode=1

%-------------------------
% Custom commands
\newcommand{\resumeItem}[1]{
  \item\small{
    {#1 \vspace{-2pt}}
  }
}

\newcommand{\resumeSubheading}[4]{
  \vspace{0pt}\item
    \begin{tabular*}{0.97\textwidth}[t]{l@{\extracolsep{\fill}}r}
      \textbf{#1} & #2 \\
      \textit{\small#3} %& \textit{\small #4} \\
    \end{tabular*} % \vspace{-7pt}
    \textit{\small #4}
    \vspace{-7pt}
}

\newcommand{\skillSubheading}[4]{
  \vspace{1pt}\item
    \begin{tabular*}{0.97\textwidth}[t]{l@{\extracolsep{\fill}}r}
      %\textbf{#1} & #2 \\
      \textit{\small#3} & \textit{\small #4} \\
    \end{tabular*}\vspace{-7pt}
}

\newcommand{\resumeSubSubheading}[2]{
    \item
    \begin{tabular*}{0.97\textwidth}{l@{\extracolsep{\fill}}r}
      \textit{\small#1} & \textit{\small #2} \\
    \end{tabular*}\vspace{-7pt}
}

\newcommand{\resumeProjectHeading}[2]{
    \item
    \begin{tabular*}{0.97\textwidth}{l@{\extracolsep{\fill}}r}
      \small#1 & #2 \\
    \end{tabular*}\vspace{-7pt}
}

\newcommand{\resumeSubItem}[1]{\resumeItem{#1}\vspace{-4pt}}

\renewcommand\labelitemii{$\vcenter{\hbox{\tiny$\bullet$}}$}

\newcommand{\resumeSubHeadingListStart}{\begin{itemize}[leftmargin=0.10in, label={}]}
\newcommand{\resumeSubHeadingListEnd}{\end{itemize}}
\newcommand{\resumeItemListStart}{\begin{itemize}[leftmargin=0.14in]}
\newcommand{\resumeItemListEnd}{\end{itemize}\vspace{-8pt}}

%-------------------------
% Page header (name + contact line). Values come from shared/profile.tex.
\newcommand{\resumeHeader}{%
  \begin{center}
    \textbf{\Huge \scshape \profileName} \\ \vspace{5pt}
    \profilePhone\ $|$
    \href{mailto:\profileEmail}{\profileEmail} $|$
    \href{https://\profileGithub}{\profileGithub} $|$
    Portfolio: \href{https://\profilePortfolio}{\profilePortfolio}
  \end{center}%
}


\begin{document}

%----------HEADING----------
\resumeHeader

%-----------EDUCATION-----------
\section{Education}
  \resumeSubHeadingListStart
    \resumeSubheading
      {Bachelor of Computer Science (Advanced)}{\small Feb 2023 -- Nov 2026}
      {Adelaide University}{}
  \resumeSubHeadingListEnd


%-----------TECHNICAL SKILLS-----------
\section{Technical Skills}
 \begin{itemize}[leftmargin=0.15in, label={}]
    \small{\item{
     \textbf{Languages}{: Python, Java, TypeScript, JavaScript, C++, C, Bash, MATLAB, GDScript} \\
     \textbf{Libraries and Frameworks}{: Node.js, React, Flask, Spring Boot, REST APIs, MySQL, MongoDB, Pandas, NumPy, Matplotlib, Scikit-learn, TensorFlow, ROS2, Nav2, Gazebo, Godot} \\
     \textbf{Tools \& Practice}{: Git, Docker, Linux, Jira, JUnit, automated testing, CI/CD} \\
    }}
 \end{itemize}


%-----------EXPERIENCE-----------
\section{Experience}
  \resumeSubHeadingListStart

  \resumeSubheading
         {\textbf{{Student Data Engineer - REDCap-KNIME Integration - Adelaide University}}}{\small Aug 2026 - Oct 2026}
         {\emph{\textbf{Tech:}} Python, KNIME, REDCap}{}
         \resumeItemListStart
           \resumeItem{Connected a research data system (REDCap) to an analytics workflow tool (KNIME) by building a custom Python extension using PyCap to pull records and metadata through the REST API into separate typed output tables.}
           \resumeItem{Implemented metadata-driven ETL and validation: decoded categorical and checkbox fields, enforced field-level validation rules, and preserved longitudinal event and repeating-instrument context.}
           \resumeItem{Built a controlled test environment of 10,000 synthetic records across 64 fields (no real participant data), enabling reproducible validation of the extraction and ETL workflow.}
         \resumeItemListEnd

  \resumeSubheading
         {\textbf{{Industry Project - Block Model Compression Algorithm - Maptek}}}{\small Aug 2025 - Oct 2025}
         {\emph{\textbf{Tech:}} C++, Bash, Jira}{}
         \resumeItemListStart
           \resumeItem{Architected a stream-oriented lossless compression system for 3D data, optimising for speed and compression ratio.}
           \resumeItem{Implemented an octree-based spatial compression algorithm, achieving ~98.8\% reduction in data size while improving processing speed by 2.6x over the uncompressed baseline, forming the foundation for subsequent team optimisations reaching 99.17\% compression and 1520x processing speed.}
           \resumeItem{Designed a testing framework that verifies correctness of compressed outputs against uncompressed input data, ensuring algorithmic correctness independent of compression strategy and enabling automated, extensible testing.}
           \resumeItem{Built a cross-platform CI pipeline with automated testing and reproducible Linux and Windows builds, integrating cross-compilation and artifact generation into the development workflow, improving team velocity by 1.5x.}
         \resumeItemListEnd

  \resumeSubheading
         {\textbf{Software Developer - Autonomous Navigation - Adelaide Rover Team}}{\small Oct 2024 - Feb 2025}
         {\emph{\textbf{Tech:}} Python, C++, ROS2, Nav2, Gazebo, Groot}{}
         \resumeItemListStart
           \resumeItem{Designed a behavior tree for the rover’s autonomous navigation pipeline using Nav2, enabling modular decision-making for tasks such as path planning, recovery behaviours, and goal execution.}
           \resumeItem{Simulated navigation behaviour using Gazebo, supporting the Adelaide Rover Team's preparation for the Australian Rover Challenge's autonomous navigation task.}
         \resumeItemListEnd

  \resumeSubHeadingListEnd


%-----------PROJECTS-----------
\section{Projects}
   \resumeSubHeadingListStart

   \resumeSubheading
         {\textbf{Daily Iftar Registration System}}{\small Feb 2026 - Mar 2026}
         {\emph{\textbf{Tech:}} TypeScript, Node.js, MongoDB, REST API, Docker}
         {\\ \small \textbf{\emph{Link:}} \href{https://aumsa.com.au/}{aumsa.com.au}}
         \resumeItemListStart
           \resumeItem{Delivered production-ready backend features under tight release cycles, building and iterating on RESTful APIs (Node.js, MongoDB) to support live event registration used by 70+ daily users.}
           \resumeItem{Architected DST-safe, timezone-aware date handling (Australia/Adelaide), eliminating off-by-one scheduling bugs and ensuring reliable, real-world consistency for day-based operations.}
           \resumeItem{Designed scalable MongoDB database schemas and API contracts to support rapid feature expansion and seamless frontend integration, maintaining stability while continuously deploying new functionality.}
         \resumeItemListEnd

   \resumeSubheading
         {\textbf{Paxos Consensus Algorithm}}{\small Sep 2025 - Oct 2025}
         {\textbf{\emph{Tech:}} Java, JUnit, Bash}
         {\\ \textbf{\emph{Link:}} \href{https://github.com/Shas-9/Paxos}{github.com/Shas-9/Paxos}}
         \resumeItemListStart
           \resumeItem{Implemented the Paxos Consensus Algorithm without use of any external libraries, enabling multiple concurrent nodes to elect a single leader despite network delays and node failures.}
           \resumeItem{Developed a TCP socket-based messaging protocol supporting Paxos phases (prepare, promise, accept, learn) with nodes dynamically acting as proposers, acceptors, and learners to achieve consensus.}
           \resumeItem{Designed a failure simulation framework introducing latency, message loss, and node crashes to validate consensus correctness under concurrent proposals and partial system failures.}
         \resumeItemListEnd

   \resumeSubheading
         {\textbf{Distributed Weather Feed Aggregator with Lamport Clocks}}{\small Aug 2025 - Sep 2025}
         {\textbf{\emph{Tech:}} Java, JUnit, REST API}
         {\\ \textbf{\emph{Link:}} \href{https://github.com/Shas-9/Weather-Aggregation-Server}{github.com/Shas-9/Weather-Aggregation-Server}}
         \resumeItemListStart
           \resumeItem{Engineered a distributed data aggregation system supporting concurrent clients and content servers.}
           \resumeItem{Implemented Lamport logical clocks to maintain causal consistency and deterministic ordering of concurrent updates across distributed components, without use of external libraries.}
           \resumeItem{Developed a multi-threaded aggregation server with thread-safe request processing, persistent storage with crash recovery and automated 30s expiry of stale data sources, ensuring system reliability under concurrent workloads.}
         \resumeItemListEnd

   \resumeSubheading
         {\textbf{SVM and Neural Network Implementation}}{\small Aug 2025 - Oct 2025}
         {\textbf{\emph{Tech:}} Python, NumPy, Sklearn}
         {\\ \textbf{\emph{Link:}} \href{https://github.com/Shas-9/SVM-and-Neural-Network-Implementation}{github.com/Shas-9/SVM-and-Neural-Network-Implementation}}
         \resumeItemListStart
           \resumeItem{Implemented a Support Vector Machine from scratch in Python by solving the dual optimization problem with quadratic programming, including RBF kernel support and hyperparameter tuning using validation data.}
           \resumeItem{Built a two-layer neural network from scratch using forward and backward propagation, implementing core modules (Linear, ReLU, Cross-Entropy) and training with mini-batch SGD on a dataset with 200 features and 8.5k samples.}
           \resumeItem{Developed an end-to-end ML training pipeline including data standardization, train/validation/test splits, hyperparameter search, and performance evaluation using accuracy metrics with Sklearn.}
         \resumeItemListEnd

   \resumeSubheading
         {\textbf{CatQuest -- 2D Maze Exploration Game}}{\small Aug 2023 - Oct 2023}
         {\emph{\textbf{Tech:}} C++, SFML}
         {\\ \textbf{\emph{Link:}} \href{https://github.com/Shas-9/Labyrinth-game}{github.com/Shas-9/Labyrinth-game}}
         \resumeItemListStart
           \resumeItem{Engineered modular game architecture, separating game state, entities, and UI for maintainability and scalability.}
           \resumeItem{Integrated enemy AI and player mechanics within the existing procedural random maze generation system, including real-time collision handling and health management.}
           \resumeItem{Designed the game UI, with seamless navigation and state transitions without affecting game loop performance.}
         \resumeItemListEnd

   \resumeSubheading
         {\textbf{Monkey Run -- Unihack 2024}}{\small Mar 2024}
         {\emph{\textbf{Tech:}} GDScript, Godot, Python, MediaPipe}
         {\\ \textbf{\emph{Link:}} \href{https://github.com/rayokamoto/unihack}{github.com/rayokamoto/unihack}}
         \resumeItemListStart
           \resumeItem{Developed an endless runner game with randomised object generation using Godot game engine with GDScript.}
           \resumeItem{Implemented input detection by tracking player movement through webcam by using MediaPipe ML framework.}
           \resumeItem{Integrated input detection system with game mechanics to deliver a cohesive gaming experience.}
         \resumeItemListEnd

     % --- Older projects, hidden from the PDF and the website. Un-comment to show. ---
     % \resumeProjectHeading
     %     {\textbf{\href{https://github.com/Shas-9/Data-visualisation}{Global GDP and Internet Users -- Research}} $|$ \emph{Python, Pandas, NumPy, Matplotlib}}{}
     %     \resumeItemListStart
     %       \resumeItem{Analysed 60 years of data sourced from Kaggle to investigate correlation between internet user growth and GDP across nations.}
     %       \resumeItem{Utilised Pandas library in Python to process data and perform necessary calculations and presented findings through visually compelling graphs generated using Matplotlib.}
     %       \resumeItem{Identified key insights to discern impact of internet penetration on economic development worldwide.}
     %     \resumeItemListEnd
     % \resumeProjectHeading
     %     {\textbf{\href{https://github.com/Shas-9/Password-manager}{Password Manager}} $|$ \emph{Python, CustomTkinter}}{}
     %     \resumeItemListStart
     %       \resumeItem{Developed a local password manager in Python, enabling convenient storage and retrieval of website and account credentials, as a practical solution for personal password management needs.}
     %       \resumeItem{Created an intuitive user interface using CustomTkinter graphics library, facilitating quick access to stored passwords through a search feature.}
     %     \resumeItemListEnd
   \resumeSubHeadingListEnd


%-----------EXTRACURRICULARS-----------
\section{Extracurriculars}
   \resumeSubHeadingListStart

   \resumeProjectHeading
         {\textbf{Adelaide University Competitive Programming League (AUCPL)}}{\small Mar 2026 -- Oct 2026}
         \resumeItemListStart
           \resumeItem{Achieved a team 2nd place finish in the Adelaide University Competitive Programming League.}
         \resumeItemListEnd

   \resumeSubheading
         {Secretary}{\small Nov 2024 -- Aug 2026}
         {ISSUA}{}
         \resumeItemListStart
           \resumeItem{Led internal communications for a 10+ member committee, producing structured meeting minutes and action item documentation to ensure organizational alignment and accountability.}
           \resumeItem{Managed university-facing logistics including room bookings, building access, and prayer room maintenance.}
           \resumeItem{Administered organizational knowledge management using Google Drive, maintaining structured repositories of meeting records, event documentation, and operational resources to support cross-functional collaboration.}
         \resumeItemListEnd

   \resumeSubheading
         {Events Officer}{\small Nov 2023 -- Nov 2024}
         {ISSUA}{}
         \resumeItemListStart
           \resumeItem{Planned and executed society events for 100+ attendees, overseeing end-to-end logistics, budget management, and stakeholder alignment to meet organizational objectives.}
           \resumeItem{Coordinated event setup, including venue preparation, infrastructure, and audio-visual systems.}
           \resumeItem{Engaged with attendees and guests during events, providing support and fostering a welcoming environment.}
         \resumeItemListEnd

   \resumeSubHeadingListEnd

%-------------------------------------------
\end{document}
